% OpenStax Calculus Volume 3 -- Section 2.2 Exercises: Vectors in Three Dimensions
% Exercises reproduced from OpenStax, licensed under CC BY 4.0.
% Source: https://openstax.org/books/calculus-volume-3/pages/2-2-vectors-in-three-dimensions
\documentclass{exercises}
\title{OpenStax Calculus Volume 3}
\subtitle{Section 2.2 Exercises: Vectors in Three Dimensions}
\sourceurl{https://openstax.org/books/calculus-volume-3/pages/2-2-vectors-in-three-dimensions}
\author{}
\date{}

\begin{document}
\maketitle

\begin{enumerate}[start=1]
\item Consider a rectangular box with one of the vertices at the origin, as shown in the following figure. If point $A(2,3,5)$ is the opposite vertex to the origin, then find \\ \textit{[Figure: This figure is the first octant of the 3-dimensional coordinate system. It has a point labeled ``A(2, 3, 5)'' drawn.]}
\begin{enumerate}[label=(\alph*)]
  \item the coordinates of the other six vertices of the box and
  \item the length of the diagonal of the box determined by the vertices $O$ and $A$.
\end{enumerate}
\item Find the coordinates of point $P$ and determine its distance to the origin. \\ \textit{[Figure: This figure is the first octant of the 3-dimensional coordinate system. It has a point drawn at (2, 1, 1). The point is labeled ``P.'']}
\end{enumerate}

For the following exercises, describe and graph the set of points that satisfies the given equation.

\begin{enumerate}[resume]
\item $(y-5)(z-6)=0$
\item $(z-2)(z-5)=0$
\item ${(y-1)}^{2}+{(z-1)}^{2}=1$
\item ${(x-2)}^{2}+{(z-5)}^{2}=4$
\item Write the equation of the plane passing through point $(1,1,1)$ that is parallel to the xy-plane.
\item Write the equation of the plane passing through point $(1,-3,2)$ that is parallel to the xz-plane.
\item Find an equation of the plane passing through points $(1,-3,-2)$, $(0,3,-2)$, and $(1,0,-2)$.
\item Find an equation of the plane passing through points $(1,9,2)$, $(1,3,6)$, and $(1,-7,8)$.
\end{enumerate}

For the following exercises, find an equation of the sphere in standard form that satisfies the given conditions.

\begin{enumerate}[resume]
\item Center $C(-1,7,4)$ and radius $4$
\item Center $C(-4,7,2)$ and radius $6$
\item Diameter $PQ$, where $P(-1,5,7)$ and $Q(-5,2,9)$
\item Diameter $PQ$, where $P(-16,-3,9)$ and $Q(-2,3,5)$
\end{enumerate}

For the following exercises, find the center and radius of the sphere with an equation in general form that is given.

\begin{enumerate}[resume]
\item $x^{2}+y^{2}+z^{2}-4z+3=0$
\item $x^{2}+y^{2}+z^{2}-6x+8y-10z+25=0$
\end{enumerate}

For the following exercises, express vector $\overset{\to}{PQ}$ with the initial point at $P$ and the terminal point at $Q$


(a) in component form and


(b) by using standard unit vectors.

\begin{enumerate}[resume]
\item $P(3,0,2)$ and $Q(-1,-1,4)$
\item $P(0,10,5)$ and $Q(1,1,-3)$
\item $P(-2,5,-8)$ and $M(1,-7,4)$, where $M$ is the midpoint of the line segment $PQ$
\item $Q(0,7,-6)$ and $M(-1,3,2)$, where $M$ is the midpoint of the line segment $PQ$
\item Find terminal point $Q$ of vector $\overset{\to}{PQ}=\langle7,-1,3\rangle$ with the initial point at $P(-2,3,5)$.
\item Find initial point $P$ of vector $\overset{\to}{PQ}=\langle-9,1,2\rangle$ with the terminal point at $Q(10,0,-1)$.
\end{enumerate}

For the following exercises, use the given vectors $\text{a}$ and $\text{b}$ to find and express the vectors $\text{a}+\text{b}$, $4\text{a}$, and $-5\text{a}+3\text{b}$ in component form.

\begin{enumerate}[resume]
\item $\text{a}=\langle-1,-2,4\rangle$, $\text{b}=\langle-5,6,-7\rangle$
\item $\text{a}=\langle3,-2,4\rangle$, $\text{b}=\langle-5,6,-9\rangle$
\item $\text{a}=-\text{k}$, $\text{b}=-\text{i}$
\item $\text{a}=\text{i}+\text{j}+\text{k}$, $\text{b}=2\text{i}-3\text{j}+2\text{k}$
\end{enumerate}

For the following exercises, vectors u and v are given. Find the magnitudes of vectors $\text{u}-\text{v}$ and $-2\text{u}$.

\begin{enumerate}[resume]
\item $\text{u}=2\text{i}+3\text{j}+4\text{k}$, $\text{v}=-\text{i}+5\text{j}-\text{k}$
\item $\text{u}=\text{i}+\text{j}$, $\text{v}=\text{j}-\text{k}$
\item $\text{u}=\langle2\cos t,-2\sin t,3\rangle$, $\text{v}=\langle0,0,3\rangle$, where $t$ is a real number.
\item $\text{u}=\langle0,1,\sinh t\rangle$, $\text{v}=\langle1,1,0\rangle$, where $t$ is a real number.
\end{enumerate}

For the following exercises, find the unit vector in the direction of the given vector $\text{a}$ and express it using standard unit vectors.

\begin{enumerate}[resume]
\item $\text{a}=3\text{i}-4\text{j}$
\item $\text{a}=\langle4,-3,6\rangle$
\item $\text{a}=\overset{\to}{PQ}$, where $P(-2,3,1)$ and $Q(0,-4,4)$
\item $\text{a}=\overset{\to}{OP}$, where $P(-1,-1,1)$
\item $\text{a}=\text{u}-\text{v}+\text{w}$, where $\text{u}=\text{i}-\text{j}-\text{k}$, $\text{v}=2\text{i}-\text{j}+\text{k}$, and $\text{w}=-\text{i}+\text{j}+3\text{k}$
\item $\text{a}=2\text{u}+\text{v}-\text{w}$, where $\text{u}=\text{i}-\text{k}$, $\text{v}=2\text{j}$, and $\text{w}=\text{i}-\text{j}$
\item Determine whether $\overset{\to}{AB}$ and $\overset{\to}{PQ}$ are equivalent vectors, where $A(1,1,1),B(3,3,3),P(1,4,5)$, and $Q(3,6,7)$.
\item Determine whether the vectors $\overset{\to}{AB}$ and $\overset{\to}{PQ}$ are equivalent, where $A(1,4,1)$, $B(-2,2,0)$, $P(2,5,7)$, and $Q(-3,2,1)$.
\end{enumerate}

For the following exercises, find vector $\text{u}$ with a magnitude that is given and satisfies the given conditions.

\begin{enumerate}[resume]
\item $\text{v}=\langle7,-1,3\rangle$, $\|\text{u}\|=10$, $\text{u}$ and $\text{v}$ have the same direction
\item $\text{v}=\langle2,4,1\rangle$, $\|\text{u}\|=15$, $\text{u}$ and $\text{v}$ have the same direction
\item $\text{v}=\langle2\sin t,2\cos t,1\rangle$, $\|\text{u}\|=2$, $\text{u}$ and $\text{v}$ have opposite directions for any $t$, where $t$ is a real number
\item $\text{v}=\langle3\sinh t,0,3\rangle$, $\|\text{u}\|=5$, $\text{u}$ and $\text{v}$ have opposite directions for any $t$, where $t$ is a real number
\item Determine a vector of magnitude $5$ in the direction of vector $\overset{\to}{AB}$, where $A(2,1,5)$ and $B(3,4,-7)$.
\item Find a vector of magnitude $2$ that points in the opposite direction than vector $\overset{\to}{AB}$, where $A(-1,-1,1)$ and $B(0,1,1)$. Express the answer in component form.
\item Consider the points $A(2,\alpha,0),B(0,1,\beta)$, and $C(1,1,\beta)$, where $\alpha$ and $\beta$ are negative real numbers. Find $\alpha$ and $\beta$ such that $\|\overset{\to}{OA}-\overset{\to}{OB}+\overset{\to}{OC}\|=\|\overset{\to}{OB}\|=4$.
\item Consider points $A(\alpha,0,0),B(0,\beta,0)$, and $C(\alpha,\beta,\beta)$, where $\alpha$ and $\beta$ are positive real numbers. Find $\alpha$ and $\beta$ such that $\|\overset{\to}{OA}+\overset{\to}{OB}\|=\sqrt{2}$ and $\|\overset{\to}{OC}\|=\sqrt{3}$.
\item Let $P(x,y,z)$ be a point situated at an equal distance from points $A(1,-1,0)$ and $B(-1,2,1)$. Show that point $P$ lies on the plane of equation $-2x+3y+z=2$.
\item Let $P(x,y,z)$ be a point situated at an equal distance from the origin and point $A(4,1,2)$. Show that the coordinates of point $P$ satisfy the equation $8x+2y+4z=21$.
\item The points $A,B$, and $C$ are collinear (in this order) if the relation $\|\overset{\to}{AB}\|+\|\overset{\to}{BC}\|=\|\overset{\to}{AC}\|$ is satisfied. Show that $A(5,3,-1)$, $B(-5,-3,1)$, and $C(-15,-9,3)$ are collinear points.
\item Show that points $A(1,0,1)$, $B(0,1,1)$, and $C(1,1,1)$ are not collinear.
\item \techrequired A force $\text{F}$ of $50\,\text{N}$ acts on a particle in the direction of the vector $\overset{\to}{OP}$, where $P(3,4,0)$.
\begin{enumerate}[label=(\alph*)]
  \item Express the force as a vector in component form.
  \item Find the angle between force $\text{F}$ and the positive direction of the x-axis. Express the answer in degrees rounded to the nearest integer.
\end{enumerate}
\item \techrequired A force $\text{F}$ of $40\,\text{N}$ acts on a box in the direction of the vector $\overset{\to}{OP}$, where $P(1,0,2)$.
\begin{enumerate}[label=(\alph*)]
  \item Express the force as a vector by using standard unit vectors.
  \item Find the angle between force $\text{F}$ and the positive direction of the x-axis.
\end{enumerate}
\item If $\text{F}$ is a force that moves an object from point $P_{1}(x_{1},y_{1},z_{1})$ to another point $P_{2}(x_{2},y_{2},z_{2})$, then the displacement vector is defined as $\text{D}=(x_{2}-x_{1})\text{i}+(y_{2}-y_{1})\text{j}+(z_{2}-z_{1})\text{k}$. A metal container is lifted $10$ m vertically by a constant force $\text{F}$. Express the displacement vector $\text{D}$ by using standard unit vectors.
\item A box is pulled $4$ yd horizontally in the x-direction by a constant force $\text{F}$. Find the displacement vector in component form.
\item The sum of the forces acting on an object is called the resultant or net force. An object is said to be in static equilibrium if the resultant force of the forces that act on it is zero. Let $\text{F}_{1}=\langle10,6,3\rangle$, $\text{F}_{2}=\langle0,4,9\rangle$, and $\text{F}_{3}=\langle10,-3,-9\rangle$ be three forces acting on a box. Find the force $\text{F}_{4}$ acting on the box such that the box is in static equilibrium. Express the answer in component form.
\item \techrequired Let $\text{F}_{k}=\langle1,k,k^{2}\rangle$, $k=1,...,n$ be $n$ forces acting on a particle, with $n\ge2$.
\begin{enumerate}[label=(\alph*)]
  \item Find the net force $\text{F}=\sum_{k=1}^{n}F_{k}$. Express the answer using standard unit vectors.
  \item Use a computer algebra system (CAS) to find n such that $\|\text{F}\|<100$.
\end{enumerate}
\item The force of gravity $\text{F}$ acting on an object is given by $\text{F}=m\text{g}$, where m is the mass of the object (expressed in kilograms) and $\text{g}$ is acceleration resulting from gravity, with $\|\text{g}\|=9.8\,\text{N/kg}$. A 2-kg disco ball hangs by a chain from the ceiling of a room.
\begin{enumerate}[label=(\alph*)]
  \item Find the force of gravity $\text{F}$ acting on the disco ball and find its magnitude.
  \item Find the force of tension $\text{T}$ in the chain and its magnitude.\\ Express the answers using standard unit vectors.
\end{enumerate}
\item A 5-kg pendant chandelier is designed such that the alabaster bowl is held by four chains of equal length, as shown in the following figure. \\ \textit{[Figure: This figure shows a light fixture hung from a ceiling, supported by 4 chains from the same point on the ceiling to four points spread evenly around the light fixture.]}
\begin{enumerate}[label=(\alph*)]
  \item Find the magnitude of the force of gravity acting on the chandelier.
  \item Find the magnitudes of the forces of tension for each of the four chains (assume chains are essentially vertical).
\end{enumerate}
\item \techrequired A 30-kg block of cement is suspended by three cables of equal length that are anchored at points $P(-2,0,0)$, $Q(1,\sqrt{3},0)$, and $R(1,-\sqrt{3},0)$. The load is located at $S(0,0,-2\sqrt{3})$, as shown in the following figure. Let $\text{F}_{1}$, $\text{F}_{2}$, and $\text{F}_{3}$ be the forces of tension resulting from the load in cables $RS,QS$, and $PS$, respectively. \\ \textit{[Figure: This figure is the 3-dimensional coordinate system. It has 4 points drawn. The first point is labeled ``P(-2, 0, 0).'' The second point is labeled ``R(1, -squareroot of 3, 0).'' The third point is labeled ``S(0, 0, -2squareroot of 3).'' The fourth point is labeled ``Q(1, squareroot of 3, 0).'' There are line segments from P to S, from R to S, and from Q to S. At point S there is a box labeled ``30 kg.'']}
\begin{enumerate}[label=(\alph*)]
  \item Find the gravitational force $\text{F}$ acting on the block of cement that counterbalances the sum $\text{F}_{1}+\text{F}_{2}+\text{F}_{3}$ of the forces of tension in the cables.
  \item Find forces $\text{F}_{1}$, $\text{F}_{2}$, and $\text{F}_{3}$. Express the answer in component form.
\end{enumerate}
\item Two soccer players are practicing for an upcoming game. One of them runs 10 m from point A to point B. She then turns left at $90^\circ$ and runs 10 m until she reaches point C. Then she kicks the ball with a speed of 10 m/sec at an upward angle of $45^\circ$ to her teammate, who is located at point A. Write the velocity of the ball in component form. \\ \textit{[Figure: This figure is the image of two soccer players. The first soccer player is at point A. The second player is at point C. There is a line segment from A to C. Ther is a vector from player C upwards labeled ``v.'' There is a vector from player A to the bottom of the image. The point at the bottom is labeled ``B.'' This vector is labeled ``10m.'' There is a vector from C to B labeled ``10m.'']}
\item Let $\text{r}(t)=\langle x(t),y(t),z(t)\rangle$ be the position vector of a particle at the time $t\in[0,T]$, where $x,y$, and $z$ are smooth functions on $[0,T]$. The instantaneous velocity of the particle at time $t$ is defined by vector $\text{v}(t)=\langle x'(t),y'(t),z'(t)\rangle$, with components that are the derivatives with respect to $t$, of the functions x, y, and z, respectively. The magnitude $\|\text{v}(t)\|$ of the instantaneous velocity vector is called the speed of the particle at time t. Vector $\text{a}(t)=\langle x''(t),y''(t),z''(t)\rangle$, with components that are the second derivatives with respect to $t$, of the functions $x,y$, and $z$, respectively, gives the acceleration of the particle at time $t$. Consider $\text{r}(t)=\langle \cos t,\sin t,2t\rangle$ the position vector of a particle at time $t\in[0,30]$, where the components of $\text{r}$ are expressed in centimeters and time is expressed in seconds.
\begin{enumerate}[label=(\alph*)]
  \item Find the instantaneous velocity, speed, and acceleration of the particle after the first second. Round your answer to two decimal places.
  \item Use a CAS to visualize the path of the particle---that is, the set of all points of coordinates $(\cos t,\sin t,2t)$, where $t\in[0,30]$.
\end{enumerate}
\item \techrequired Let $\text{r}(t)=\langle t,2t^{2},4t^{2}\rangle$ be the position vector of a particle at time $t$ (in seconds), where $t\in[0,10]$ (here the components of $\text{r}$ are expressed in centimeters).
\begin{enumerate}[label=(\alph*)]
  \item Find the instantaneous velocity, speed, and acceleration of the particle after the first two seconds. Round your answer to two decimal places.
  \item Use a CAS to visualize the path of the particle defined by the points $(t,2t^{2},4t^{2})$, where $t\in[0,10]$.
\end{enumerate}
\end{enumerate}

\end{document}
